\section{Policy Experiments and Counterfactuals}\label{sec:policy}

The model transforms fiscal policy from a question of aggregate size into a
question of distributional design. This section quantifies that claim through
four policy experiments that span the key margins on which governments make
real decisions: (i) targeted vs.\ untargeted public spending; (ii) austerity
and the asymmetric welfare cost of fiscal consolidation; (iii) automatic
stabilizers under countercyclical incidence; (iv) developing-economy
applications where $\lambda$ is large and public-good quality is low.

\subsection{Experiment 1: targeted vs.\ untargeted spending}
\label{subsec:targeted-untargeted}

\paragraph{Design.}
A 1\% fiscal expansion (quarterly, persistence $\rho_{g}=0.8$) is allocated in
two ways. Under the \emph{untargeted} baseline, $g$ rises uniformly across
households, matching the standard BoE scenario of broad-based public-sector
procurement. Under the \emph{targeted} counterfactual, $g$ rises only in
complementary categories consumed primarily by the bottom half of the wealth
distribution: public health, public transport, education, and housing support.

\paragraph{Result.}
Targeted fiscal expansion generates an impact multiplier of $\mu=2.08$, compared
to $\mu=1.68$ for untargeted spending -- a 24\% increase in the multiplier
from re-allocating the \emph{same} total fiscal outlay. The welfare gain is
concentrated in the bottom quintile (a 1.8\% consumption-equivalent gain
relative to the no-policy counterfactual) but the top quintile is also
marginally better off (+0.2\%) because the aggregate output expansion raises
their labor income.

\paragraph{Mechanism.}
The amplification comes entirely from the $\xi$ channel: targeted spending
hits households with larger $|\theta^{H}|$, raising the effective cross-
elasticity $\epsilon^{H}_{cg}$ at the aggregate level. In the language of
Corollary~\ref{cor:het}, the distributional dominance condition is
strengthened because the weight $\lambda$ effectively rises -- from the 0.35
share of \HtM{} households overall, to the 0.65 share of bottom-half
households exposed to the targeted spending.

\subsection{Experiment 2: austerity and asymmetric welfare}
\label{subsec:austerity}

\paragraph{Design.}
A 1\% fiscal \emph{contraction} in public services (health, transit, education)
under two scenarios: (a) separable benchmark ($\theta^{H}=\theta^{S}=0$) and
(b) heterogeneous internalization ($\theta^{H}=-0.76$, $\theta^{S}=+0.29$). The
consolidation is implemented as a persistent reduction in spending on
complementary categories.

\paragraph{Result.}
The output contraction is $1.8\%$ in the separable benchmark and $2.6\%$ under
heterogeneous internalization -- a \emph{$44\%$ larger contractionary effect}.
Even more importantly, the welfare cost is sharply asymmetric: the bottom
quintile suffers a 2.2\% consumption-equivalent loss, while the top quintile
suffers only 0.3\%.

\paragraph{Policy implication.}
Fiscal consolidations implemented by cutting public services have welfare costs
that are \emph{larger} than conventional multiplier estimates suggest and are
concentrated among the households least able to bear them. A model that
suppresses the $c$--$g$ complementarity channel will systematically
under-estimate both the aggregate and distributional costs of austerity. This
result directly informs the recent IMF debate on the role of fiscal
consolidation in stabilization programs \citep{Blanchard2019} and suggests that
fiscal rules should incorporate differential treatment of spending categories.

\subsection{Experiment 3: automatic stabilizers}
\label{subsec:automatic}

\paragraph{Design.}
We compare two stabilizer configurations against the no-stabilizer baseline:
(a) a uniform transfer that pays $\Delta$ dollars to every household during a
recession; (b) a targeted public-service maintenance program that preserves
the level of complementary public goods (health, education, transit) during
the downturn. Both stabilizers are equal in total fiscal cost.

\paragraph{Result.}
The targeted stabilizer delivers an output stabilization ratio of 1.32
(measured as the reduction in output volatility relative to the no-stabilizer
baseline), compared to 1.18 for the uniform transfer. More striking, the
targeted stabilizer delivers a welfare gain to the bottom quintile that is
3.1 times larger than the uniform transfer. The mechanism is again the $\xi$
channel: maintaining public services during downturns preserves the
marginal-utility support of private consumption for the most exposed
households, which dampens the consumption response to the shock.

\paragraph{Policy implication.}
Automatic stabilizers operating through public-service maintenance are an
under-appreciated tool. The conventional macroeconomic case for stabilizers is
based on output smoothing, which our results reinforce. But the case is
strengthened by the \emph{distributional} dimension that only emerges when
$c$--$g$ complementarity is explicit.

\subsection{Experiment 4: developing-economy applications}
\label{subsec:developing}

\paragraph{Design.}
We re-calibrate the model to a representative Sub-Saharan African economy
following \citet{Francois2023}, with the following changes relative to the US
baseline: (a) $\lambda=0.70$ (higher \HtM{} share); (b) $\theta^{H}=-1.20$
(stronger \HtM{} complementarity -- \citet{Gethin2023}'s ``triple curse'' of
regressive incidence); (c) $\bar g/\bar y=0.10$ (lower public-goods share in
GDP); (d) a Taylor rule with a higher weight on output stabilization.

\paragraph{Result.}
The impact multiplier in the developing-economy calibration is $\mu=2.64$,
compared to $\mu=1.68$ in the US baseline. Two forces drive the amplification:
a larger $\lambda$ mechanically increases the NK Cross, and a larger
$|\theta^{H}|$ raises the direct fiscal channel $\xi$.

\paragraph{Policy implication.}
The quantitative prediction that fiscal multipliers are substantially larger
in low-income countries is consistent with the IMF's revised multiplier
estimates for developing economies \citep{IMF2020} and provides a
micro-founded mechanism for the empirical pattern. It also strengthens the
case for fiscal expansion as a poverty-reduction tool in these economies: the
distributional dominance condition holds strongly, so fiscal expansion targeted
at complementary public services delivers both output amplification and
progressive distributional effects simultaneously.

\subsection{Summary of policy implications}

Table~\ref{tab:policy-summary} summarizes the multipliers and welfare effects
across the four experiments.

\begin{table}[ht]
\centering\small
\caption{Policy experiments: impact multipliers and welfare effects.}
\label{tab:policy-summary}
\begin{threeparttable}
\begin{tabular}{@{}lrrr@{}}
\toprule
 & {Impact $\mu$} & {Bottom Q welfare \%} & {Top Q welfare \%} \\
\midrule
US untargeted (baseline)            & 1.680 & +0.9 & +0.1 \\
US targeted to complements          & 2.080 & +1.8 & +0.2 \\
Austerity (separable benchmark)     & -1.800 & -1.1 & -0.2 \\
Austerity (heterogeneous internalization) & -2.600 & -2.2 & -0.3 \\
Automatic stabilizer (uniform)      & \multicolumn{1}{c}{\text{1.18 var.}} & +0.5 & +0.1 \\
Automatic stabilizer (targeted)     & \multicolumn{1}{c}{\text{1.32 var.}} & +1.5 & +0.1 \\
Developing economy (baseline)       & 2.640 & +2.7 & +0.4 \\
\bottomrule
\end{tabular}
\begin{tablenotes}\footnotesize
\item Welfare effects measured as consumption-equivalent lifetime change
relative to the no-policy counterfactual. The stabilizer rows report the
output-volatility-reduction ratio rather than an impact multiplier. All
experiments use the baseline calibration of Table~\ref{tab:calibration}
modified as described in the text.
\end{tablenotes}
\end{threeparttable}
\end{table}

The four experiments converge on a single take-home: fiscal policy design
matters as much as fiscal policy size. Allocating a fixed fiscal envelope
toward categories used by constrained households generates larger aggregate
multipliers and more progressive welfare effects simultaneously. Cutting those
categories has both larger and more regressive contractionary effects. The
distributional incidence of public goods is therefore a first-order
macroeconomic consideration -- not merely a redistribution question.
